\documentclass[10pt,a4paper]{amsart}
\usepackage[margin=19mm]{geometry}
\usepackage{amsmath,amssymb,mathtools}
\usepackage[scr=rsfs,cal=euler]{mathalfa}
\usepackage{xfrac}
\usepackage[unicode,hidelinks,bookmarks=false]{hyperref}
\newcommand{\ie}{i.e.\ }
\newcommand{\cf}{cf.\ }
\newcommand{\resp}{respectively\ }
\newcommand{\Subsection}{\subsection}
\providecommand{\nobreakdash}{\nobreak\hbox{-}}
\setlength{\emergencystretch}{2em}
\multlinegap0pt
\usepackage{bm}
\usepackage{bbm}
\usepackage{dsfont}
\usepackage{xspace}
\usepackage{cancel}
\usepackage{mathtools}
\usepackage{amsthm}
\makeatletter
\@ifundefined{theo}{\newtheorem{theo}{Theo}}{}
\@ifundefined{lemm}{\newtheorem{lemm}[theo]{Lemma}}{}
\makeatother
\newcommand{\cO}{\mathcal{O}} 
\newcommand{\dcorn}{$\mathsf{d}$-cornet}
\newcommand{\dtrum}{$\mathsf{d}$-trumpet}
\newcommand{\out}{\mathrm{deg_{out}}}
\newcommand{\rE}{\mathrm{E}} 
\newcommand{\rV}{\mathrm{V}} 
\renewcommand{\rm}{\mathrm{m}}
\title[Blossoming bijection for bipartite maps: a new approach via ]{Blossoming bijection for bipartite maps: a new approach via orientations and applications to the Ising model}
\author{Marie Albenque, Laurent M\'enard, Nicolas Tokka}
\date{Source: arXiv:2511.03680v2 (CC BY 4.0)}
\begin{document}
\maketitle

\noindent\emph{Source and licence.} Blossoming bijection for bipartite maps: a new approach via orientations and applications to the Ising model. Version arXiv:2511.03680v2 (CC BY 4.0), distributed under the Creative Commons Attribution 4.0 International licence, as recorded in the arXiv licence metadata for this exact version and shown on the abstract page https://arxiv.org/abs/2511.03680v2. Full rights notice: licenses/arxiv-2511.03680-CC-BY-4.0.txt.\par\smallskip

\noindent\textit{Editorial scope.} Counted unit: the Lemm whose source label is \texttt{lemm: Characterization tightness} in the article (source0.tex lines 857--871, source label \texttt{lemm: Characterization tightness}), delivered together with the article's own complete original proof of it (source0.tex lines 873--944, one \texttt{proof} environment of 72 lines). No mathematics is written, repaired or re-derived: statement and proof are the article's own text line for line. Required context retained uncounted: none. Every definition, display and label of those lines is retained unchanged. Omitted from this package and not redistributed: 1--856, 872--872, 945--1995. Nothing inside a retained excerpt is omitted or altered: every retained body is delivered exactly as the article's lines are written, so no omission receipt of any kind accompanies this package. Citations: the retained excerpts carry no citation command at all, so no bibliography entry of the article is transcribed and no reference is redistributed. Reference adaptation: none was needed and none was made; every reference inside the delivered unit resolves inside it, so no unresolved reference or citation is produced. Printed slips kept literal: no printed word, symbol, hypothesis or number of the counted statement or of its proof was altered. Layout and class: the article's own class is \texttt{amsart} with packages including amsthm, amssymb, amsfonts, amsmath, amscd, tikz-cd, tikz, mathrsfs, subcaption, graphicx, vmargin, hyperref, xspace, array; the wrapper is the lane's pinned \texttt{amsart} editorial wrapper at 10pt on A4, which restates the theorem-like environments the retained text needs and adds the article's own macro definitions that the retained text needs (\texttt{\textbackslash cO}, \texttt{\textbackslash dcorn}, \texttt{\textbackslash dtrum}, \texttt{\textbackslash out}, \texttt{\textbackslash rE}, \texttt{\textbackslash rV}, \texttt{\textbackslash rm}), with extra wrapper packages (bm, bbm, dsfont, xspace, cancel, mathtools). The delivered numbering is the unit's own. Assets: no retained span contains a figure environment, a graphics-inclusion command, a PDF-page inclusion, a table or a TikZ picture; no image file is copied into this package, the package declares no asset dependency and no image file is redistributed with it. Licence: version arXiv:2511.03680v2 of this article is distributed under the Creative Commons Attribution 4.0 International licence, as recorded in the arXiv licence metadata for this exact version and shown on the abstract page https://arxiv.org/abs/2511.03680v2. A full rights notice is delivered at licenses/arxiv-2511.03680-CC-BY-4.0.txt. Characters: every retained excerpt is pure ASCII, so the delivered engine, XeLaTeX with its default Latin Modern text font, reports no missing character.
\begin{lemm}\label{lemm: Characterization tightness}
Fix $d\geq k >0$.  Let $(\rm,\tau)$ be a pair with $\Delta(\rm)\leq d$, and $\deg(\tau)=k$.

If $\tau$ is black, then: 
\[
	(m,\tau)\text{ is a \dtrum}\quad\Leftrightarrow\quad \alpha_{d,k}^-\text{ is feasible on }(\rm,\tau),
\]
and, in that case, any $\alpha_{d,k}^-$-orientation is accessible.

If $\tau$ is white, then: 
\[
	(m,\tau)\text{ is a \dcorn}\quad\Leftrightarrow\quad \alpha_{d,k}^+\text{ is feasible and quasi-accessible on }(\rm,\tau),
\]
where quasi-accessible means that for any $v\neq \tau$, there exists a forward path from $v$ to $\rho$. 
\end{lemm}

\begin{proof}
Assume that $\tau$ is black, and that $\alpha_{d,k}^-$ is feasible on $(\rm,\tau)$. Let $(R,S)$ be a black cut of $\rm$, and let $w\geq 1$ denote its weight. We aim to prove that $w\geq k$.  Let $d_\bullet$ and $d_\circ$ denote the sum of the degrees of the black and white vertices in $R$, respectively. First, note that since $(R,S)$ is black, we have that $d_\circ = d_\bullet+w$. Then, since  $\alpha_{d,k}^-$ is $(d+1)$-fractional, we have: 
\begin{equation}\label{eq:ori_cut}
	\sum_{v\in R}{\alpha_{d,k}^-(v)} \geq \left(d+1\right) \vert\mathrm{E}_R\vert,
\end{equation}
where $ \vert\mathrm{E}_R\vert$ is the number of edges with both endpoints in $R$. Moreover, on one hand, by definition of $\alpha_{d,k}^-$-orientations, we have: 
\[
	\sum_{v\in R}{\alpha_{d,k}^-(v)} = d\, d_\bullet + d_\circ -k.
\]
On the other hand, since $(R,S)$ is black, we have: 
\[
\left(d+1\right) \vert\mathrm{E}_R\vert	=\left(d+1\right) d_\bullet = d\, d_\bullet + d_\circ - w.
\]
By~\eqref{eq:ori_cut}, it follows that $w\geq k$, and hence that $(\rm,\tau)$ is a \dtrum. \\

Reciprocally, assume that $(\rm,\tau)$ is a \dtrum. To define an $\alpha_{d,k}^-$-orientation on $(\rm,\tau)$, we apply the max-flow min-cut theorem to the following representation of $(\rm,\tau)$ as a flow network.
The \emph{flow network} associated to $\rm$, and denoted by $\tilde{\rm}$, is the directed plane map obtained by replacing each edge $e$ of $\rm$ with $2$ directed copies: one toward its black endpoint, denoted by $e_{\circ\rightarrow\bullet}$, and the other toward its white endpoint, denoted by $e_{\bullet\rightarrow\circ}$. Moreover the \emph{capacities} of the edges are set to be: 
\[
\textrm{cap}(e_{\circ\rightarrow\bullet})\coloneqq 1\quad \text{ and } \quad \textrm{cap}(e_{\bullet\rightarrow\circ})\coloneqq d, \text{ for any }e\in \rE(\rm).
\]
Finally, the \emph{source} of $\tilde\rm$ is $\rho$ and its \emph{sink} is $\tau$. Recall that a \emph{flow} is a function $F:\mathrm{E}(\tilde\rm)\rightarrow \mathbb{Z}_{\geq 0}$ that must satisfy the following conditions:
\begin{itemize}
    \item Capacity constraint: $\forall \vec e\in \mathrm{E}(\tilde \rm)$, $F(\vec e) \leq \textrm{cap}(\vec e)$, 
    \item Conservation of flows: $\forall v \in V\setminus \{ \rho, \tau\}, \sum_{\vec{e}\in\mathrm{E}_{v \rightarrow}}{F (\vec{e})} = \sum_{\vec{e}\in\mathrm{E}_{\rightarrow v}}{F(\vec{e})},$
    \item Source/sink constraint: $\sum_{\vec{e}\in\mathrm{E}_{\rightarrow \rho}}{F(\vec{e})} = \sum_{\vec{e}\in\mathrm{E}_{\tau \rightarrow}}{F(\vec{e})}= 0$,
\end{itemize}
where $\mathrm{E}_{v \rightarrow}$ and $\mathrm{E}_{\rightarrow v}$ denote the set of edges in $\tilde\rm$ directed from $v$ and toward $v$, respectively.
Moreover, the value $|F|$ of the flow is defined by: 
\[
|F| \coloneqq  \sum_{\vec{e}\in\mathrm{E}_{\rho \rightarrow}}{F(\vec{e})} = \sum_{\vec{e}\in\mathrm{E}_{\rightarrow \tau}}{F(\vec{e})}.
\]


The \emph{capacity} of a cut $(R,S)$ of $\tilde\rm$ is the sum of the capacity of the edges of its cut-set that are directed toward a vertex in $S$ if the source is in $R$, or toward a vertex in $R$ if the source is in $S$. We write $\underline{\mathrm{cut}}$ for the minimal capacity of the cuts of $\tilde\rm$, and prove that $\underline{\mathrm{cut}} = k$. Observe that if $(R,S)$ is black (resp. white), then its capacity is equal to its weight (resp. to $d$ times its weight).

Since $\tau$ is black, the capacity of the cut defined by the trivial partition $(V\backslash\{\tau\},{\tau})$ is equal to $k$. Then, let $(R,S)$ be any other cut of $\tilde{\rm}$. Either it corresponds to a black cut of $\rm$, and since $\rm$ is tight, it follows that its capacity is at least $k$. Or, there is at least an edge $\vec{e}$ in the cut-set that is directed toward a white vertex in $S$. Since the capacity of $\vec{e}$ is equal to $d\geq k$, the capacity of $(R,S)$ is at least $k$, which implies that $\underline{\mathrm{cut}} = k$. 

The max-flow min-cut theorem then ensures the existence of a flow $F^\star$ with value $k$, that we use to define an explicit orientation $\cO$ on $\rm$. For every edge $e=\left\{h_\bullet,h_\circ\right\}\in\mathrm{E}(\rm)$, we set:
\begin{equation*}
\begin{aligned}
	\cO(h_\circ) &\coloneqq  1 + F^\star(e_{\bullet\rightarrow\circ}) - F^\star(e_{\circ\rightarrow\bullet})\\
	\cO(h_\bullet) &\coloneqq  d - F^\star(e_{\bullet\rightarrow\circ}) + F^\star(e_{\circ\rightarrow\bullet}).
\end{aligned}
\end{equation*}
First, it follows from the \emph{capacity constraints} that $-1\leq F^\star(e_{\bullet\rightarrow\circ}) - F^\star(e_{\circ\rightarrow\bullet})\leq d$, for every edge $e\in\mathrm{E}(\rm)$. Hence, $\cO$ is a well defined $(d+1)$-fractional orientation. Then, it follows from the conservation of flows at each vertex and from $\vert F^\star\vert = k$ that $\out(v)=\alpha_{d,k}^-(v)$, for every $v\in\mathrm{V}(\rm)$. Thus, $\cO$ is an $\alpha_{d,k}^-$-orientation on $(\rm,\rho,\tau)$. \

To conclude, note that $\cO$ is accessible.
Indeed, for any subset $S\subset \mathrm{V}(\rm)\setminus\{\rho\}$, we have:
\begin{equation*}
\sum_{v\in S}{\alpha_{d,k}^-(v)} \geq \sum_{v\in S}{\alpha_d(v)} > \left(d+1\right) \vert\mathrm{E}_S\vert,
\end{equation*}
where the last inequality follows from the accessibility of the $\alpha_d$-orientations on the bipartite map $\rm$.\\


The case of $\alpha_{d,k}^+$-orientation is similar. The only subtlety is to deal with the uniqueness of the minimal cut (for the converse implication) and with the quasi-accessibility (for  the direct implication). For the uniqueness of the minimal cut, using the quasi-accessibility, we can replace~\eqref{eq:ori_cut} by: 
\[
\sum_{v\in S}\alpha_{d,k}^+(v)  > \left(d+1\right) \vert\mathrm{E}_S\vert,
\]
where $(R,S)$ is any white cut different from the trivial one. The uniqueness of the minimal cut follows from the same arguments as in the previous case. 

For the quasi-accessibility, we define the flow network and the associated orientation in the same way, except that we exchange the role of the source (which becomes $\tau$) and of the sink (which becomes $\rho$), while the definition of $\cO$ remains the same. The same arguments as in the previous case show that the capacity of any cut is at least $k$, with equality only for the trivial cut.

Let $(R,S)\neq (\rV(\rm)\backslash\{\tau\},\{\tau\})$ be a black cut, and let $w$ denote its weight. The following holds by the conservation of flows at each vertex and from $\vert F^\star\vert = k$:
\begin{equation*}
	\sum_{v\in S}\alpha_{d,k}^+(v) = \sum_{v\in S}\sum_{h\sim v} \cO(h) = \left(d+1\right) \vert\mathrm{E}_S\vert + w - \vert F^\star\vert.
\end{equation*}
Since, $(\rm,\rho,\tau)$ is a \dtrum, one has $w>\vert F^\star\vert$, so that:
\begin{equation*}
	\sum_{v\in S}\alpha_{d,k}^+(v) > \left(d+1\right) \vert\mathrm{E}_S\vert.
\end{equation*}
The quasi-accessibility follows. 
\end{proof}

\end{document}
